% ────────────────────────────────────────────────────────
\section{Limites e Continuidade}

% ────────────────────────────────────────────────────────
\subsection{Enquadramento Teórico}

\begin{defbox}[title={Limite em $\mathbb{R}^2$}]
  Diz-se que $\displaystyle\lim_{(x,y)\to(a,b)}f(x,y) = L$ se para todo o
  $\varepsilon>0$ existe $\delta>0$ tal que
  \[
    0 < \sqrt{(x-a)^2+(y-b)^2} < \delta \implies |f(x,y)-L|<\varepsilon.
  \]
\end{defbox}

\begin{notebox}[title={Estratégia: Provar que um Limite Não Existe}]
  Ao contrário de $\mathbb{R}$, em $\mathbb{R}^2$ o limite tem de ser o mesmo
  ao longo de \emph{todos} os caminhos. Uma estratégia comum para mostrar que
  um limite \emph{não existe} consiste em exibir dois caminhos que dão valores
  diferentes.
\end{notebox}

\begin{defbox}[title={Continuidade em $\mathbb{R}^2$}]
  $f$ é \emph{contínua} em $(a,b)$ se
  $\displaystyle\lim_{(x,y)\to(a,b)}f(x,y)=f(a,b)$.
\end{defbox}

% ────────────────────────────────────────────────────────
\subsection{Exemplos Resolvidos}

\begin{exbox}{4 --- Limite Inexistente}
  Mostre que $\displaystyle\lim_{(x,y)\to(0,0)}\frac{xy}{x^2+y^2}$ não existe.
\end{exbox}
\begin{solbox}
  Pelo caminho $y=0$: $f(x,0)=0\to 0$.\\
  Pelo caminho $y=x$: $f(x,x)=\dfrac{x^2}{2x^2}=\dfrac{1}{2}\to\dfrac{1}{2}$.\\
  Como os limites são diferentes, o limite não existe.
\end{solbox}
